\section{Sliding Window Minimum}
\label{sec:cses-sw-minimum}

\problemstatement{Given an array of $n$ integers and a window size $k$,
output the minimum of every contiguous window of size $k$.}

\begin{patternbox}
Minimum has no cheap inverse, so a running value fails. The right tool is a
\textbf{monotonic deque} that keeps candidate minima in increasing order:
the front is always the current window's minimum, and each element is
pushed and popped at most once, giving $O(n)$ overall.
\end{patternbox}

\subsection*{The monotonic deque}

Maintain a deque of \emph{indices} whose values are strictly increasing
from front to back. When a new element arrives, pop from the back every
index whose value is $\ge$ the new value -- those can never be the minimum
again while the newer, smaller element is in the window. Then push the new
index. Pop from the front any index that has slid out of the window
(index $\le i-k$). The front index's value is the window minimum.

\begin{ideabox}[Discard Dominated Candidates]
An older element that is larger than a newer one is ``dominated'': as long
as the newer element is present, the older can never be the minimum, so it
is removed from the back forever. Each index enters and leaves the deque
once, which is why the total work is linear despite the inner popping.
\end{ideabox}

\subsection*{Algorithm}

\begin{enumerate}
  \item For each index $i$: pop back while its value $\ge a[i]$; push $i$.
  \item Pop front while it is $\le i-k$ (out of window).
  \item Once $i\ge k-1$, output $a[\text{front}]$.
\end{enumerate}

\subsection*{Dry run}

\dryrun{$a=[2,1,3,4,1]$, $k=2$.}

\begin{itemize}
  \item $i{=}0$: deque $[0]$. $i{=}1$: $a[1]{=}1\le a[0]$, pop $0$, push
        $1$; window min $a[1]=1$.
  \item $i{=}2$: push $2$ (deque $[1,2]$, but $1$ now out at $i-k=0$?
        front $1>0$ stays); min $a[1]=1$. Then front $1\le2-2=0$? no.
  \item Continuing yields minima $1,1,3,1$ for the four windows.
\end{itemize}

\subsection*{C++ implementation}

\begin{lstlisting}
#include <bits/stdc++.h>
using namespace std;

int main() {
    int n, k;
    cin >> n >> k;
    vector<int> a(n);
    for (int& x : a) cin >> x;

    deque<int> dq;                            // indices, values increasing
    for (int i = 0; i < n; i++) {
        while (!dq.empty() && a[dq.back()] >= a[i]) dq.pop_back();
        dq.push_back(i);
        if (dq.front() <= i - k) dq.pop_front();   // drop out-of-window
        if (i >= k - 1) cout << a[dq.front()] << " ";
    }
    cout << "\n";
    return 0;
}
\end{lstlisting}

\begin{complexitybox}
\textbf{Time Complexity.} $O(n)$ -- each index is pushed and popped at most
once. \textbf{Space Complexity.} $O(k)$ for the deque.
\end{complexitybox}

\begin{takeawaybox}
The monotonic deque is the canonical answer to sliding min/max: keep only
non-dominated candidates in monotonic order, front is the answer. It is the
single most reused sliding-window structure and appears again as the
maximum (advertisement) variant.
\end{takeawaybox}
